\section{Related Work}
\label{sec:related}

Our work sits at the intersection of visual appearance robustness, counterfactual \textsc{VLM} diagnosis, and 3D embodied simulation. Table~\ref{tab:related_work} summarizes how our protocol differs from prior literature.

\begin{table*}[!t]
\centering\scriptsize
% Prior robustness/bias benchmarks compared with ConvertAsset
\begin{tabular*}{\textwidth}{@{\extracolsep{\fill}}p{3.1cm}p{3.0cm}p{3.2cm}p{2.9cm}p{3.7cm}@{}}
\toprule
Work & Stimulus / Domain & Perturbation Mechanism & Primary Metrics & Key Distinction vs.\ Ours \\
\midrule
\citet{Geirhos2019Texture}; \citet{Burgert2025NotTexture} & 2D ImageNet categories & Neural style transfer or post-hoc cue suppression & Top-1 acc., shape vs.\ texture ratio & We keep geometry fixed via 3D shader conversion on indoor scenes. \\
\citet{Hendrycks2019Benchmarking} & ImageNet-C / -P corruptions & Additive 2D pixel noise, blur, weather, digital & mCE, flip prob.\ on frame sequences & Our converted renders stay clean (1024$\times$768) with zero sensor corruption. \\
\citet{Zhang2024ImageNetD} & Diffusion synthetic objects & Unconstrained generative backgrounds/textures & Classification accuracy drop & Deterministic USD shader mapping avoids generative shape drift. \\
\citet{Gavrikov2025Talk}; \citet{Eppel2025Shape} & 2D/3D cue-conflict \& LAS\&T & Prompt steering; multi-factor crop variation & Category / material accuracy & Fixed pointing prompt measuring pixel-exact localization turnover. \\
\citet{Tong2024EyesWideShut}; \citet{Sar2026RealityCheck} & MMVP; Reality Check VQA & CLIP-blind pairs; 2D edit / swap counterfactuals & VQA accuracy across bias axes & Native 3D BSDF swaps paired with $\mathrm{CIE}~L^*a^*b^*$ salience matching. \\
\citet{Wang2026ConstantEye} & \emph{navdream} driving & Generative pixel-aligned style transfer & Driving trajectories \& adapter stability & Indoor SimReady MDL$\to$Preview conversion across 4 VLM families + risk gating. \\
\citet{Chaubey2020UAP} & Adversarial robustness survey & Gradient-optimized universal pixel noise & Attack success / fooling rate & Unoptimized engineering export induces bidirectional gains and losses. \\
\midrule
\textbf{Ours (\textsc{ConvertAsset})} & \textbf{30 USD scenes + KuJiaLe} & \textbf{3D MDL $\to$ Preview + 1,440 mask CFs} & \textbf{Paired turnover, L1--L4 floor, AUROC} & \textbf{4 VLM families (6 ckpts), 2-family CF (+3.54\%), zero-VLM risk gate.} \\
\bottomrule
\end{tabular*}

\caption{Comparison between prior appearance-robustness, counterfactual-bias, and adversarial studies and our \textsc{ConvertAsset} material-conversion framework.}
\label{tab:related_work}
\end{table*}

\paragraph{Texture, shape, and appearance robustness}
Classic texture--shape cue-conflict experiments established that ImageNet-trained convolutional networks rely heavily on local surface texture compared to human observers \citep{Geirhos2019Texture}, while \textsc{ImageNet-C} and \textsc{ImageNet-P} systematized robustness to additive 2D corruptions and frame perturbations \citep{Hendrycks2019Benchmarking}. More recent controlled cue-suppression studies refine that picture, showing that local shape remains competitive when confounds are carefully controlled \citep{Burgert2025NotTexture}. In multimodal models, \citet{Gavrikov2025Talk} show that \textsc{VLM}s lean more toward global shape than isolated vision encoders yet remain steerable toward surface cues via prompting, while \citet{Eppel2025Shape} benchmark 2D and 3D material/texture recognition across massive automatically extracted crops. Generative diffusion suites such as \textsc{ImageNet-D} \citep{Zhang2024ImageNetD} expose severe accuracy drops under synthesized backgrounds and materials, and \emph{The Constant Eye} \citep{Wang2026ConstantEye} isolates autonomous-driving style shifts via pixel-aligned generative transfer. Unlike 2D pixel corruptions or diffusion re-synthesis---which either degrade image sharpness or risk subtle silhouette drift---3D shader translation from \textsc{MDL} to \usdps{} alters surface optics within clean, geometry-locked ray-traced scenes encountered directly in robotics pipelines.

\paragraph{Counterfactual VLM diagnostics and adversarial perturbations}
Diagnostic benchmarks increasingly employ paired interventions to separate visual understanding from language priors. \citet{Tong2024EyesWideShut} identify \textsc{CLIP}-blind pairs where semantic encoders conflate visually distinct stimuli, while \emph{VLM Reality Check} \citep{Sar2026RealityCheck} uses minimal 2D edits, swaps, and silhouettes to audit cognitive biases in visual question answering. Methodologically, our counterfactuals combine natural 3D shader conversion with pixel-exact silhouette compositing and $\mathrm{CIE}~L^*a^*b^*$ photometric matching across two distinct 8B \textsc{VLM} families to disentangle local luminance/color salience from non-photometric \textsc{BSDF} structure. Importantly, unlike universal adversarial perturbations \citep{Chaubey2020UAP}---which optimize imperceptible worst-case noise using model gradients to maximize error---industrial shader lowering uses zero model feedback and induces \emph{bidirectional} effects (removing distracting specular glare on some targets while eroding surface cues on others).

\paragraph{Spatial grounding interfaces and embodied simulation}
Modern open-weight \textsc{VLM}s expose diverse referential pointing and grounding conventions across distinct backbones: dynamic-resolution \textsc{Qwen2.5-VL} and \textsc{Qwen3-VL} \citep{Bai2025Qwen,Bai2025Qwen3}, \textsc{InternVL3}/\textsc{InternVL3.5} \citep{Zhu2025InternVL,Wang2025InternVL35}, specialist pointing architectures \textsc{Molmo}/\textsc{Molmo2} built on SigLIP-2 \citep{Deitke2024Molmo,Clark2026Molmo2}, open \textsc{CLIP}-Mistral architectures \textsc{LLaVA-NeXT} \citep{Liu2024LLaVANeXT}, as well as earlier grounding models \citep{Chen2023Shikra,You2024Ferret,Google2026Gemma4,Agrawal2024Pixtral}. Because prompt sensitivity and coordinate scaling can easily masquerade as perceptual errors \citep{Zhao2021Calibrate}, every model in our study passes an explicit 24-input pilot qualification gate before 30-scene evaluation (\ref{app:readout}). Beyond static images, Vision-and-Language Navigation (\textsc{VLN}) \citep{Anderson2018Vision} and interactive \textsc{USD} ecosystems such as \textsc{GRUtopia} \citep{Wang2024GRUtopia} and \textsc{DualVLN} \citep{Wei2026Ground} evaluate sequential closed-loop policies, complementing traditional domain randomization \citep{Tobin2017Domain,Tremblay2018Training} by testing whether standard cross-renderer shader lowering preserves task-level agent decisions.
